%==============================
\chapter{Conclusion}\label{cap:conclusao}
%==============================

This work investigated the capability of five Large Language Models—ChatGPT, Gemini, Claude, Meta AI, and DeepSeek—to convert sequential C programs into parallel OpenMP implementations. 
The evaluation was conducted on six programs from the PAMPAR benchmark, covering a range of parallelization scenarios from embarrassingly parallel reductions to wavefront algorithms with complex data dependencies. 
Each model was given the same sequential source code and a single standardized prompt, and the generated programs were evaluated for compilation success, output correctness, execution speedup at 11 threads, and the quality of the parallelization strategy.

%==============================
\section{Summary of Findings}
\label{sec:conclusion_summary}

\textbf{RQ1 (Correctness).} All 30 generated programs—five models across six programs—compiled without errors and produced output identical to the sequential oracle. 
The 100\% correctness rate across a diverse set of programs, including programs with complex data dependencies such as the wavefront Turing Ring simulation, demonstrates that current frontier LLMs have sufficient understanding of C programming and the OpenMP model to produce functionally valid parallel code from a single zero-shot prompt.

\textbf{RQ2 (Speedup).} Consistently (averaging over DP, NI, MM, OE, and DFT, with Turing Ring excluded uniformly for all models), three of the five models closely matched the PAMPAR expert-written reference (5.11$\times$): Gemini and Claude at 5.17$\times$ each, and DeepSeek at 5.15$\times$. ChatGPT was lower (4.79$\times$ average), driven primarily by a suboptimal matrix multiplication implementation (2.76$\times$). 
Meta AI achieved the lowest average (4.22$\times$), driven by a serializing Numerical Integration implementation (0.48$\times$). 
On individual programs, LLMs frequently matched or exceeded the reference: all five models surpassed the reference speedup for Dot Product (reference 4.59$\times$; LLMs 5.35--5.43$\times$), and all five achieved near-reference speedups for DFT (reference 6.15$\times$; LLMs 5.99--6.10$\times$).

\textbf{RQ3 (Failure Characterization).} Two qualitatively different failure modes were identified. 
First, Meta AI applied \texttt{\#pragma omp ordered} to the Numerical Integration accumulation loop, serializing a computation that requires no ordering constraints. 
The correct solution — a \texttt{reduction(+:)} clause — was used by all other models; this failure shows that a model may select a technically valid synchronization primitive that is inappropriate for the actual data dependency structure. 
Second, ChatGPT and Meta AI both handled the Turing Ring wavefront correctly, with dependency ordering.
But their inner parallel loops iterate over the full $N \times N$ grid per diagonal step rather than bounding the iteration to valid cells on each diagonal. 
That produces $O(N^3)$ iterations where $O(N^2)$ suffice, and creates one parallel region per diagonal step, leading to substantial threading spawn overhead. 
Notably, DeepSeek independently derived the $O(N^2)$ bounded iteration as part of its solution, achieving a 34.75$\times$ speedup that reflects an algorithmic improvement rather than parallelism alone.

\textbf{RQ4 (Model Comparison).} The five models were indistinguishable in terms of correctness. In terms of performance, Gemini, Claude, and DeepSeek were the strongest performers, consistently close to the expert reference across the programs. 
ChatGPT and Meta AI exhibited program-specific weaknesses that significantly reduced their overall averages: ChatGPT in the matrix multiplication application and Meta AI in numerical integration. 
The observed differences are not attributable to a single consistent cause; rather, each model's shortfalls, arising from distinct decisions in specific programs, suggest that no single parallelization failure mode dominates current LLMs.

%==============================
\section{Limitations}
\label{sec:conclusion_limitations}

Several limitations of this study should be acknowledged. 
First, the evaluation uses a single prompt per program without iterative refinement. In practice, a programmer would likely refine an LLM's output through follow-up prompts, which could resolve some of the observed failures. 
Second, the benchmark is limited to six programs from PAMPAR; different programs or larger scales might reveal different strengths and weaknesses. 
Third, the models were accessed through web interfaces, which may use different underlying model versions than the APIs, and model versions change over time, limiting the reproducibility of exact results. 
Fourth, the study uses a fixed 11-thread configuration on a single hardware platform; the relative performance of different parallelization strategies may differ on platforms with more threads or different memory hierarchies. 
Finally, a single generation attempt per program per model was used; sampling multiple generations and selecting the best could yield higher performance.

%==============================
\section{Future Work}
\label{sec:conclusion_future}

Several directions extend the present work. 
Iterative prompting experiments—in which the model is shown its own speedup output and asked to improve it—could assess whether LLMs can self-correct their parallelization strategies. 
Extending the evaluation to additional PAMPAR programs (particularly JA, GL, and GS), other parallel programming interfaces (MPI, CUDA, OpenACC), and compiled-and-tested feedback loops would provide a more comprehensive picture. 
A study of whether the failure modes observed here—incorrect synchronization of primitive selection and unbounded iteration—consistently appears across different program inputs and hardware configurations would support more general conclusions. 
Finally, comparing the LLM-generated code against automatic parallelizing compilers on the same programs would place the LLM approach in the broader context of the automatic parallelization literature.

%==============================
%\section{About AI use}
%We use ...


